\documentclass[11pt]{scrartcl}

\usepackage[stixtwo]{fontsetup}
\usepackage{amsmath}

\begin{document}
Ha a sorozat korlátos és monoton, akkor konvergens, és határértéke a sorozat
legkisebb felső vagy legnagyobb alsó korlátja. A lehetséges $A$ határérték
kielégíti az $A = 4 - \frac{3}{A}$ egyenletet, vagyis $A = 1$ vagy $A = 3$.

Belátjuk, hogy $1 < a_n < 3$: $1 < a_1 < 3$, és ha $1 < a_n < 3$, akkor
$1 > \frac{1}{a_n} > \frac{1}{3}$, tehát $-3 < -\frac{3}{a_n} < -1$, így
\[
	1 = 4 - 3 < a_{n + 1} = 4 - \frac{3}{a_n} < 4 - 1 = 3.
\]

$a_2 = \frac{5}{2} > a_1$, és ha $a_n < a_{n + 1}$, akkor
$\frac{1}{a_n} > \frac{1}{a_{n + 1}}$, és
$-\frac{3}{a_n} < -\frac{3}{a_{n + 1}}$, tehát
\[
	a_{n + 1} = 4 - \frac{3}{a_n} < 4 - \frac{3}{a_{n + 1}} = a_{n + 2},
\]
vagyis a sorozat monoton nő. Mivel $a_n \ge 2$, így a határérték csak $3$
lehet.

A sorozat tehát monoton és korlátos, így konvergens, a határértéke $3$.
\end{document}
